%! Polynomial Long and Synthetic Division — Unit 3, Lesson 3.5 — Lesson Plan
% Section order per LESSON_SHAPE.md §5. CED 1.11.B (EK 1.11.B.1, the division algorithm; the
% slant-asymptote EK 1.11.B.2 is skipped — no rational functions in this unit) with CED 1.5.A.1
% (a zero <=> a linear factor), plus the course's additions: synthetic division by x - c, the
% Remainder Theorem, and the Factor Theorem. 3.4 factored by pattern; 3.5 factors by division when
% no pattern applies; 3.6 runs the chain backwards (zeros -> polynomial).
% Crux: the MISSING POWER — notes 11 vs 12 (x^3 - 7x + 6 / (x - 2) without and with the 0x^2
% placeholder), settled by 13 (p(2) = 0) and explained by 14; You Do 16 and 17 carry it.
% Second trap: the SIGN OF c — notes 4 and 10; You Do 15 and 18.
\documentclass[10pt]{article}
\usepackage{precalculus-article}
\usepackage{precalculus-boxes}
\usepackage{graphicx}   % -article does NOT load graphicx; needed for screenshots
\graphicspath{{images/}}
\newcommand{\TallMath}[1]{$\displaystyle #1\rule[-1.4em]{0pt}{3.2em}$}
\newcommand{\CourseName}{Precalculus}
\newcommand{\MeetingLength}{60 minutes}
\newcommand{\UnitNumberName}{Unit 3: Polynomial Functions \quad}
\newcommand{\LessonNumberName}{Lesson 3.5: Polynomial Long and Synthetic Division}

\begin{document}
\begin{center}
    {\Large\bfseries \CourseName \\
    {\small\bfseries \UnitNumberName \LessonNumberName}}
\end{center}

\begin{tcolorbox}[colback=lilac, colframe=plum, boxrule=0.9pt, arc=2mm,
    left=2mm, right=2mm, top=1.5mm, bottom=1.5mm]
    \textbf{Primary Objective:} Students will divide one polynomial by another with long
    division and write the result as dividend $=$ divisor $\cdot$ quotient $+$ remainder, with
    the remainder lower in degree than the divisor; divide by $x - c$ with synthetic division,
    writing a $0$ for every missing power; use the \textit{Remainder Theorem} to read $p(c)$ as
    the remainder; and use the \textit{Factor Theorem} --- $(x - c)$ is a factor exactly when
    $p(c) = 0$ --- to divide out one known zero of a cubic and find all of its zeros.
    \hfill\textit{\small (CED Topics 1.11 and 1.5 --- Skills 1.B, 3.B)}
\end{tcolorbox}

\boxguard[22]
\begin{skillbox}[Lesson Flow --- Gradual Release (60 minutes)]{lilac}
\small
\renewcommand{\arraystretch}{1.25}
\begin{tabularx}{\linewidth}{@{} l c X @{}}
\textbf{Phase} & \textbf{Min} & \textbf{What students are doing, and where it lives} \\
\midrule
Warm-Up                       & 5  & Silent, individual spiral review --- warm-up sheet. \\
Hook                          & 3  & Whole-class discussion --- hook box atop the guided notes. \\
\textbf{I Do} --- model       & 10 & Watch and annotate while the teacher thinks aloud --- the LONG DIVISION and SYNTHETIC rows: definitions, the two side-by-side displays, problems 1 and 3. \\
\textbf{We Do} --- guided     & 20 & Fill the notes together, every problem cold-called --- problems 2 and 4--6, then the REMAINDER THEOREM and FACTOR THEOREM rows (7--14; the teacher works 7 and 9). \\
\textbf{You Do} --- independent & 17 & Work alone while the teacher circulates --- the You Do row, problems 15--18. \\
Debrief                       & 5  & Whole-class share-out, the headline, and the DeltaMath assignment. \\
\end{tabularx}

\smallskip
\textbf{Two procedures and two theorems share one block,} so the I Do models both procedures
on the same cubic --- problem 3 is problem 1 done the short way, and students see the same
numbers come out. There is no activity sheet to raid when the guided phase overruns. If the
clock slips, cut \textbf{problem 2} (division by $x^2 + 1$; DeltaMath's long-division topics
carry it) and \textbf{problem 5} first, and protect the FACTOR THEOREM row: problems 11--13 are
the lesson, and You Do problems 16 and 17 cannot be done without them. Never cut the You Do
below 14 minutes --- circulating it is the only formative read of the period now that there is
no exit ticket.
\end{skillbox}

\renewcommand{\arraystretch}{1.6}
\begin{skillbox}[Priority Ideas \& Skills]{goldbox}
\begin{tabularx}{\linewidth}{@{}X X@{}}
\textbf{Priority Skills} &
\textbf{Key Understandings} \\
\midrule
\textbf{AP Skill 1.B} --- Express functions, equations, or expressions in analytically
equivalent forms. &
\begin{itemize}[leftmargin=1.2em,itemsep=2pt,topsep=0pt]
  \item Dividing $f(x)$ by $g(x)$ produces a quotient $q(x)$ and a remainder $r(x)$ with
        $f(x) = g(x)\,q(x) + r(x)$ and $\deg r < \deg g$ --- the same shape as
        $157 = 12 \cdot 13 + 1$. (1.11.B.1)
  \item Long division runs divide $\to$ multiply $\to$ subtract $\to$ bring down until the
        remainder's degree drops below the divisor's.
  \item When the divisor is $x - c$, synthetic division does the same arithmetic on the
        coefficients alone: it \emph{adds} multiples of $c$ where long division subtracts
        multiples of $-c$. It works only for a divisor of the form $x - c$.
\end{itemize} \\
\textbf{AP Skill 3.B} --- Describe the assumptions and limitations of a method or model. &
\begin{itemize}[leftmargin=1.2em,itemsep=2pt,topsep=0pt]
  \item \textbf{Remainder Theorem:} dividing $p(x)$ by $x - c$ leaves remainder $p(c)$,
        because $p(c) = (c - c)\,q(c) + r = r$.
  \item \textbf{Factor Theorem:} $(x - c)$ is a factor of $p(x)$ exactly when $p(c) = 0$ ---
        remainder $0$, factor, zero, and $x$-intercept are one fact. (1.5.A.1)
  \item Knowing one zero of a cubic, divide its factor out; the quadratic quotient factors by
        3.4's patterns and gives the other zeros.
\end{itemize} \\
\textbf{The distinction the lesson is built on.} &
\begin{itemize}[leftmargin=1.2em,itemsep=2pt,topsep=0pt]
  \item \textbf{Target misconception --- the missing power:} dividing $x^3 - 7x + 6$ with no
        $0x^2$ placeholder slides every number one column left, so the division answers a
        different question (it divides $x^2 - 7x + 6$). Notes problems 11 and 12 run the two
        divisions side by side: remainders $-4$ and $0$; problem 13 settles it with
        $p(2) = 0$. You Do 16 and 17 carry it.
  \item \textbf{Second trap --- the sign of $c$:} dividing by $x + 2$ uses $c = -2$, because
        $x + 2 = x - (-2)$. Notes problems 4 and 10; You Do 15 and 18.
\end{itemize} \\
\end{tabularx}
\end{skillbox}

\begin{skillbox}[{Vocabulary, Concepts, \& Theorems}]{greenbox}
\renewcommand{\arraystretch}{1.4}
\begin{tabularx}{\linewidth}{@{} >{\raggedright\arraybackslash\leavevmode\bfseries\color{plum}}p{0.24\linewidth} X @{}}
Division algorithm & For polynomials $f$ and $g \ne 0$: $f(x) = g(x)\,q(x) + r(x)$, where $q$ is the quotient and the remainder $r$ has lower degree than $g$ (or is $0$). \\
Quotient and remainder & What a division produces: how many times the divisor goes in, and what is left over. \\
Placeholder & A $0$ coefficient written for a missing power, so every term keeps its own column. \\
Synthetic division & A shortcut for dividing by $x - c$ that works on the coefficients alone: bring down, multiply by $c$, add, repeat. \\
Remainder Theorem & When $p(x)$ is divided by $x - c$, the remainder is $p(c)$. \\
Factor Theorem & $(x - c)$ is a factor of $p(x)$ if and only if $p(c) = 0$. \\
\end{tabularx}
\end{skillbox}

\begin{skillbox}[Activate Prior Knowledge \& Spiral Review (5 min)]{lilac}
    Please see attached warm-up (spiral review).
    Skills reviewed: numerical long division, written as dividend $=$ divisor $\cdot$
    quotient $+$ remainder --- the exact shape of today's division algorithm;
    evaluating a cubic at a negative input with the input in parentheses --- the
    substitution road of the Remainder Theorem;
    factoring a quadratic trinomial and naming its zeros (3.4) --- where every quotient in the
    Factor Theorem row ends up;
    and combining like terms when one power is missing from a polynomial --- the placeholder
    idea, met first where nothing goes wrong.
\end{skillbox}

\begin{skillbox}[Hook (3 min)]{lilac}
    \textbf{No pattern fits.} Write $p(x) = x^3 - 7x + 6$ on the board.
    \begin{enumerate}[itemsep=2pt]
        \item Ask the class to run the 3.4 checklist on it: a common factor? grouping (there
              are only three terms)? a sum or difference of cubes? None applies.
        \item Read the hook box's two voices: Maya, ``There is no way to tell whether $x - 2$
              is a factor,'' and Jordan, ``Just plug in 2.'' Take two or three answers to the
              question: \textit{how could plugging in one number tell you anything about a
              factor?}
        \item Let someone compute $p(2) = 8 - 14 + 6 = 0$ and connect it to 3.3: a zero at 2
              means $(x - 2)$ is a factor. Then leave the real question hanging --- if it is a
              factor, what is the \textit{other} factor, and how do we get it?
    \end{enumerate}
    \textit{Discussion prompt:} the same polynomial returns as the crux in the FACTOR THEOREM
    row, where two careful-looking divisions disagree about it.
\end{skillbox}

\begin{skillbox}[Lesson --- I Do / We Do / You Do]{lilac}
\begin{multicols}{2}

\textbf{I DO --- Part 1: Long Division and Its Shortcut} (10 min)

\textit{Notes rows LONG DIVISION and SYNTHETIC: read each definition, mark up both displays,
work problems 1 and 3. Teacher works; students watch and annotate. Nothing is cold-called in
this phase.}
\begin{itemize}
  \item Start with the numbers in the row 1 display: $157 \div 12$ is 13 remainder 1, so
        $157 = 12 \cdot 13 + 1$. Then the polynomial beside it, step for step: $x + 3$
        remainder 1. Fill the two label boxes yourself --- $(x + 3)$ and $1$.
  \item Work \textbf{problem 1} in its frame, saying the four steps aloud every cycle:
        divide, multiply, subtract (``change every sign''), bring down. Write the result in
        division-algorithm form before moving on.
  \item Read the SYNTHETIC sentences, then put the two displays side by side: the long
        division of $(x^3 - 2x^2 - 5x + 6) \div (x - 3)$ and its synthetic twin. Point at
        $1,\ 1,\ -2$ in both. Fill ``coefficients'' and ``remainder.''
  \item Work \textbf{problem 3} --- problem 1 the short way --- and let the class see that
        $2,\ 1,\ -3 \mid 0$ is the quotient they already found.
\end{itemize}

\textbf{WE DO --- Part 2: Practice and the Sign of $c$} (6 min)

\textit{Problems 2 and 4--6. Class fills the blanks together; cold-call every one.}
\begin{itemize}
  \item \textbf{Problem 2} divides by $x^2 + 1$: the division stops when the remainder,
        $-2x + 1$, has degree 1 --- less than the divisor's 2.
  \item \textbf{Problem 4} is the sign trap: ask for $c$ \emph{before} anyone writes a number.
        $x + 2 = x - (-2)$, so $c = -2$. \textbf{Problem 5} is routine.
  \item \textbf{Problem 6} names the limit of the shortcut: $x^2 + 1$ is long division only.
\end{itemize}

\textbf{WE DO --- Part 3: The Remainder Theorem} (5 min)

\textit{Notes row REMAINDER THEOREM. Class fills the blanks together; cold-call every one.
Teacher works problem 7; the class works 8.}
\begin{itemize}
  \item Read the theorem, then the two roads in the display: substitution gives 13, the
        synthetic row gives 13. Fill ``$0$'' and ``remainder'' and say why: at $x = 3$ the
        factor $(x - 3)$ is worth nothing.
  \item Work \textbf{problem 7} both ways. \textbf{Problem 8} uses the theorem with no
        division at all, then turns it around: a remainder of 0 means $p(c) = 0$.
\end{itemize}

\columnbreak

\textbf{WE DO --- Part 4: The Factor Theorem and the Missing Power} (9 min)

\textit{Notes row FACTOR THEOREM --- the spine of the lesson. Class fills the blanks together;
cold-call every one. Teacher works problem 9; the class works 10--14.}
\begin{itemize}
  \item Read the theorem and walk the four-link chain in the display. On the graph of the
        hook's $p(x) = x^3 - 7x + 6$, fill $p(2) = 0$ and ``is.''
  \item Work \textbf{problem 9} from problem 1's quotient: $(x - 2)(2x + 3)(x - 1)$, zeros
        $2,\ 1,\ -\tfrac{3}{2}$. \textbf{Problem 10} gives a zero, $-3$, and asks for the
        factor --- the sign trap again: $x + 3$, $c = -3$.
  \item \textbf{Problems 11 and 12 are the crux.} Half the room finishes Ana's row
        ($1,\ -7,\ 6$), half finishes Ben's ($1,\ 0,\ -7,\ 6$). Put both bottom rows on the
        board: remainder $-4$ against remainder $0$. Do not resolve it yet.
  \item \textbf{Problem 13} resolves it with the Remainder Theorem: $p(2) = 0$, so Ben.
        \textbf{Problem 14} explains it: Ana's answer multiplies back to $x^2 - 7x + 6$ ---
        she divided a different polynomial. Read the lilac box aloud; give one minute for the
        ``in your own words'' line.
\end{itemize}

\textbf{YOU DO --- Part 5: Divide \& Factor} (17 min)

\textit{Notes row DIVIDE \& FACTOR, problems 15--18 --- the whole independent phase. Students
work alone; circulate and say nothing a neighbor could say instead.}
\begin{itemize}
  \item \textbf{Problem 15} divides by $x + 2$ (the sign of $c$) and writes the division
        algorithm; \textbf{16} finds a remainder two ways on $x^4 - 3x^2 + 5$, which needs
        \emph{two} placeholders --- the crux transferred.
  \item \textbf{Problem 17} proves $x - 4$ is a factor of $x^3 - 13x - 12$ (a missing $x^2$
        again) and finds all three zeros.
  \item \textbf{Problem 18} is the finisher's question: find $k$ so that $x + 2$ is a factor.
        It needs the Factor Theorem run backwards --- set $p(-2) = 0$ and solve.
\end{itemize}
\end{multicols}
\end{skillbox}

\begin{skillbox}[Explicit Instruction: Synthetic Division]{lilac}
\begin{minipage}[t]{0.48\linewidth}
\textbf{Steps}
\begin{enumerate}[itemsep=2pt]
  \item Check the divisor is $x - c$. Name $c$: $x - 3 \to 3$; $x + 2 \to -2$.
  \item Write every coefficient of the dividend, highest power first --- a $0$ for each
        missing power.
  \item Bring the first coefficient down.
  \item Multiply it by $c$; write the product under the next coefficient; add.
  \item Repeat to the last column. The last number is the remainder ($= p(c)$); the others are
        the quotient's coefficients, one degree lower.
\end{enumerate}
\end{minipage}\hfill
\begin{minipage}[t]{0.48\linewidth}
\textbf{Worked example:} $(x^3 - 7x + 6) \div (x - 2)$

\smallskip
\begin{center}
{\setlength{\tabcolsep}{4pt}\begin{tabular}{c|cccc}
 $2$ & $1$ & $0$ & $-7$ & $6$ \\
     &     & $2$ & $4$  & $-6$ \\ \hline
     & $1$ & $2$ & $-3$ & $\mathbf{0}$ \\
\end{tabular}}
\end{center}

\smallskip
Quotient $x^2 + 2x - 3 = (x + 3)(x - 1)$, remainder $0$. So
$x^3 - 7x + 6 = (x - 2)(x + 3)(x - 1)$, with zeros $2,\ -3,\ 1$. Without the $0$, the row
$1,\ -7,\ 6$ gives remainder $-4$ --- wrong, since $p(2) = 0$.
\end{minipage}
\end{skillbox}

\begin{skillbox}[Active Monitoring]{redbox}
Circulate during the \textbf{You Do} phase --- the seventeen minutes where students are working
problems 15--18 without you --- and check. The crux is \textbf{problems 11--13} (the missing
$0x^2$: remainder $-4$ against $0$, settled by $p(2) = 0$); its You Do transfers are
\textbf{problem 16} (two missing powers in $x^4 - 3x^2 + 5$) and \textbf{problem 17} (the
missing $x^2$ in $x^3 - 13x - 12$).
\begin{itemize}
  \item \textbf{Common error (the crux):} a coefficient row with no $0$ for the missing power.
        Do not say ``you forgot the zero.'' Ask: ``Read me your row with the power each number
        belongs to.'' Then: ``Does your remainder match $p(c)$?''
  \item \textbf{Common error (the second trap):} $c = 2$ for the divisor $x + 2$ in problem 15.
        Cold-call: ``What number makes $x + 2$ equal zero?''
  \item \textbf{Common error:} in long division, subtracting only the first term of the
        product --- $-(2x^3 - 4x^2)$ written but $-4x^2$ never sign-changed. Point at the
        parentheses in problem 1.
  \item \textbf{Common error:} reading the bottom row as a quotient of the \emph{same} degree,
        so $1,\ 4,\ 3$ becomes $x^3 + 4x^2 + 3x$. Ask: ``What degree is the quotient when you
        divide a cubic by $x - c$?''
  \item \textbf{Common error:} stopping at the quotient in problem 17 and listing only $4$.
        Ask: ``How many zeros can a cubic have? Where are the other two?''
\end{itemize}
\end{skillbox}

\begin{skillbox}[Differentiation --- During You Do (17 min)]{redbox}
There is no separate activity sheet: students work the notes' You Do row (problems 15--18)
alone while you circulate. Differentiate by \textit{where you stand}, not by handing out
different paper.

\begin{itemize}
  \item \textbf{Support.} Sit with a struggling student on \textbf{problem 15} and have them
        say the checklist in the You Do box aloud before writing anything: every power present
        (yes --- no placeholder needed), and $c = -2$. Then let them do the arithmetic; the
        frame does the rest.
  \item \textbf{On level.} \textbf{Problems 16 and 17} are the crux check. Do not fill in a
        placeholder for anyone; ask them to compare their remainder with $p(c)$ and walk away.
        A mismatch is the student's own signal that a column slid.
  \item \textbf{Extend.} \textbf{Problem 18} is the finisher's question: the unknown sits
        \emph{inside} the polynomial, so the Factor Theorem has to be run backwards ---
        $p(-2) = 0$, then solve for $k$. A finisher who gets $k = 3$ can be asked to divide
        $x^3 + 3x^2 - x - 6$ by $x + 2$ and confirm the remainder is $0$. Their answer
        presents in the Debrief.
\end{itemize}
\end{skillbox}

\boxguard[20]
\begin{skillbox}[Debrief (5 min)]{redbox}
Whole class, teacher-led, packets still open. Three moves, in order:
\begin{enumerate}[itemsep=3pt]
  \item \textbf{Share out the You Do} (3 min) --- one answer per problem, not a full review.
        Problem 15: $(x + 2)(2x^2 - x - 2) + 9$, with $c = -2$ said aloud. Problem 16: the
        remainder $9$, both ways, and the coefficient row $1,\ 0,\ -3,\ 0,\ 5$. Problem 17: the
        zeros $4,\ -1,\ -3$. Problem 18: $k = 3$ --- if a student solved it, they present it,
        because DeltaMath will not ask it this way.
  \item \textbf{Name the headline} (1 min) --- say it, then make the class say it back:
        \textit{dividing by $x - c$ leaves the remainder $p(c)$, so a remainder of zero means
        $x - c$ is a factor and $c$ is a zero --- and every power gets its own column, even the
        missing ones.} Point at Ana's and Ben's bottom rows while you say it.
  \item \textbf{Point forward} (1 min) --- the DeltaMath assignment, \textit{Lesson 3.5}, and
        its due date: write it on the board and have students copy it into row 4 of the cover.
        Then 3.6 runs today's chain backwards: start from the zeros and build the polynomial.
\end{enumerate}
\textbf{Do not} re-teach long division here or take new questions. If the share-out shows the
placeholder did not land, it is the 3.6 warm-up's job, not this minute's.
\end{skillbox}

\begin{skillbox}[Reinforcement \& Extension]{goldbox}
\textbf{Homework --- DeltaMath, graded.} Assign as \textit{Lesson 3.5}, in this order:
\begin{enumerate}[itemsep=1pt,topsep=2pt]
  \item \textbf{Polynomial Long Division (Level 1)} (3 problems) --- matches notes problem 1.
  \item \textbf{Polynomial Long Division (Level 2)} (2 problems) --- matches notes problem 2.
  \item \textbf{Synthetic Division (Level 1)} (3 problems) --- matches notes problems 3--5
        and 15.
  \item \textbf{Synthetic Division (Level 2)} (3 problems) --- matches notes problems 12 and
        16; missing powers and the sign of $c$.
  \item \textbf{Remainder Theorem (Level 1)} (3 problems) --- matches notes problems 7, 8,
        13, and 16.
  \item \textbf{Factor Theorem} (3 problems) --- matches notes problems 9, 13, and 18.
  \item \textbf{Synthetic Division to Find Zeros (Lev 1)} (3 problems) \textit{(diagnostic)}
        --- matches notes problems 10 and 17.
\end{enumerate}
DeltaMath has no topic aimed at the missing-power error itself: its items do include missing
terms, but they never show two divisions that disagree, so the reason rides on You Do problems
16 and 17 (and AP Practice item 5).
\textit{Leave out:} Polynomial Long Division (Level 3)--(Level 5) and Synthetic Division
(Level 3) (beyond the divisors taught), Polynomial Division Box Method (Level 1)--(Level 4) (a
third method, not taught), Find the Dividend of Polynomial Division, Remainder / Factor
Theorem (Level 2) and Synthetic Division to Find Zeros (Lev 2) (they lean on finding the first
zero unaided), Rational Root Theorem and Integer Root Theorem (not taught), and Find the
Equation from Roots (3.6).

\textbf{AP Practice --- extra credit, scored 1 to 5 (2 pages).} Optional, last in the packet.
Section I is five AP-style multiple-choice items (four options) whose distractors are the
lesson's errors: $c = +2$ for $x + 2$, the skipped placeholder, or the constant term taken as
the remainder (1); a remainder whose degree reaches the divisor's (2); $c = +1$, a quotient that
keeps the dividend's degree, or coefficients copied without dividing (3); a zero's sign flipped
or an output read as a zero (4); and the setup with the placeholder dropped, $c$ flipped, or
both (5). Section II is one free-response question on a table of a startup's monthly profit:
the Factor Theorem from $P(1) = 0$ and a synthetic division (a), every break-even month in
context (b), and an analyst's claimed remainder judged with the Remainder Theorem and
interpreted (c) --- the finisher's question.

\textbf{Preview:} Next lesson (3.6, Roots and Polynomials) runs today's chain backwards: given
the zeros, write the factors, multiply them out, and land on the polynomial in standard form.
\end{skillbox}

\begin{teachernote}[Warm-Up]
Five minutes, silent start, no calculators. Item 1 is the whole lesson in numbers --- a
student who cannot write $157 = 12 \cdot 13 + 1$ will not see why $f = g\,q + r$ is the natural
form, so name it in the I Do. Item 2 is the Remainder Theorem's first road: watch for
$-2^3$ written as $8$, which is the same sign slip that turns $c = -2$ into $c = 2$. Item 3 is
where every Factor Theorem problem ends (3.4). Item 4 quietly plants the placeholder:
$x^3 + 4x - 1$ has no $x^2$ term, and a student who lines up $4x$ under $-2x^2$ has made the
crux error in a place where it is easy to see --- point back to it in Part 4.
\end{teachernote}

\begin{teachernote}[Guided Notes]
The notes are one Main Ideas / Notes table, four instruction rows and a You Do row; in each
instruction row you read the printed sentences, mark up the display, and work the first
problem of the grid (1, 3, 7, 9), and the class works the rest with the pen in their hand,
every problem cold-called. Every division sits in a pre-drawn frame --- the long-division
frames print the divisor, the dividend and the subtraction rules; the synthetic frames print
only the L, so students must choose $c$ and write the coefficient row themselves. That choice
is where both of today's errors live, so never write a coefficient row for a student.

LONG DIVISION opens on numbers on purpose: $157 \div 12$ and $(x^2 + 5x + 7) \div (x + 2)$ both
come out ``3 remainder 1'' in their own language. SYNTHETIC is taught as \emph{the same
division}, not a new trick: the side-by-side display and problem 3 (problem 1 redone) are the
proof. The REMAINDER THEOREM row is short; its job is to give students a check they can run on
any division, which the crux then needs.

FACTOR THEOREM is the spine. Split the room for problems 11 and 12 so that both answers
exist in the room at once, and let the disagreement sit for a moment before problem 13 settles
it with $p(2) = 0$. Problem 14 is the explanation students remember: Ana divided
$x^2 - 7x + 6$. The You Do row (15--18) is the entire independent phase and carries seventeen
minutes. If the clock slips, cut problem 18 and let the Debrief carry it.
\end{teachernote}

\begin{teachernote}[Debrief]
Five minutes, and they are at risk: the FACTOR THEOREM row runs long when the crux
discussion goes well. Borrow from problem 2 or problem 5 in the We Do instead, never from the
You Do. The share-out is one answer per You Do problem so that a student who stalled on 15
still hears the coefficient row $1,\ 0,\ -3,\ 0,\ 5$ for problem 16 --- the shape of the
diagnostic DeltaMath topic, which they will otherwise meet alone at home. The point-forward
minute is the only time the DeltaMath due date is said aloud; do not drop it.
\end{teachernote}

\begin{teachernote}[AP Practice]
Extra credit, optional, scored 1 to 5, and not a substitute for the DeltaMath assignment --- say
so when you hand out the packet. Its 1-to-5 score is added as extra-credit points.
Multiple-choice answers are \textbf{C, B, D, A, D}. Item 5 is the one to glance at: (B) dropped
the placeholder for the missing $x^3$, (C) used $c = 2$ for $x + 2$, and (A) did both --- the
lesson's two errors, one per option. In item 1, (B) $18$ is the skipped-placeholder answer and
(A) $-2$ is $p(2)$. In the free response, accept a correct quotient found by long division in
(a) if the Factor Theorem reason is there; in (c), credit any answer that names $P(4) = -3$ as
the remainder and reads it as a \$3000 loss in month 4.
\end{teachernote}

\begin{teachernote}[Homework]
DeltaMath, graded, assigned as \textit{Lesson 3.5} and due on the date you set.
\textbf{Synthetic Division to Find Zeros (Lev 1)} is the diagnostic: check its item-level
results before 3.6. A wrong zero with the right magnitude and the wrong sign is the sign of $c$;
a wrong quotient on a polynomial with a missing term is the placeholder; a correct quotient
with only one zero listed means the student stopped before factoring --- the Factor Theorem
chain did not land. DeltaMath cannot show two divisions that disagree, nor ask \emph{why} a
remainder must equal $p(c)$: that lives in notes problems 13--14 and You Do 16, so read those
on the packet when you collect it. \textbf{Factor Theorem} and the diagnostic predict 3.6, where
zeros become factors again in the other direction; if either is weak, reuse a one-line item like
problem 10 in the 3.6 warm-up rather than re-teaching.
\end{teachernote}
\end{document}
